\subsection{Usecase Thanh toán}

\begin{itemize}

    \item \textbf{viewTransactionHistory(): void (Class MemberCustomer)}
    \begin{itemize}
        \item Mục đích: Kích hoạt luồng xem lịch sử giao dịch cá nhân của khách hàng có tài khoản. (UC 3.1)

        \item Mô tả logic: Phương thức này truy xuất danh sách các đối tượng Transaction liên kết với MemberCustomer hiện tại từ cơ sở dữ liệu. Sau khi lấy được dữ liệu, nó điều khiển giao diện (View/UI) hiển thị danh sách này ra màn hình, bao gồm các thông tin theo Use-case yêu cầu (Thời gian, Số tiền, Trạng thái thanh toán) và cho phép người dùng thực hiện thao tác sắp xếp.
    \end{itemize}

    \item \textbf{payInvoice(): void (Class MemberCustomer)}
    \begin{itemize}     
        \item Mục đích: Xử lý thao tác thanh toán một hóa đơn điện tử thông qua cổng thanh toán. (UC 3.2)

        \item Mô tả logic: Phương thức này tiếp nhận lệnh ``Chuyển đến BKPay'' từ người dùng và gọi đến API của cổng thanh toán BKPay. Nó chờ Webhook trả về và xử lý kết quả.      
    \end{itemize}

    \item \textbf{generateInvoice(): void (Class Invoice)}
    \begin{itemize}
        \item Mục đích: Khởi tạo thông tin chi tiết cho một hóa đơn mới. (UC 3.2)

        \item Mô tả logic: Dựa vào phí gửi xe đã được tính toán từ các phiên gửi xe (ParkingSession), hàm này khởi tạo đối tượng \texttt{Invoice}, gán \texttt{invoiceID}, thiết lập số tiền \texttt{amount}, ngày tạo \texttt{createdDate}, và đặc biệt là hạn thanh toán \texttt{dueDate}. Dữ liệu này là tiền đề để hiển thị danh sách hóa đơn cho \texttt{MemberCustomer} thanh toán.
    \end{itemize}

    \item \textbf{getDetail(): string (Class Transaction)}
    \begin{itemize}
        \item Mục đích: Đóng gói và trả về chuỗi thông tin chi tiết của một giao dịch.
        \item Mô tả logic: Phương thức này định dạng các thuộc tính của Transaction thành một chuỗi có cấu trúc dễ đọc. Chuỗi này được sử dụng để hiển thị trong danh sách lịch sử giao dịch của khách hàng.
    \end{itemize}

    \item \textbf{updateStatus(): void (Class Transaction)}
    \begin{itemize}
        \item Mục đích: Thay đổi trạng thái thanh toán của giao dịch.

        \item Mô tả logic: Truy xuất các thuộc tính của lớp \texttt{Transaction} (như \texttt{transactionID}, \texttt{amount}, \texttt{timestamp}, \texttt{paymentStatus}, \texttt{paymentMethod}) và định dạng chúng thành một chuỗi (String) hoặc chuỗi JSON có cấu trúc.
        \begin{itemize}
            \item Trong U3.2: Được gọi sau khi \texttt{payInvoice()} nhận kết quả trả về từ BKPay để đổi trạng thái từ ``Unpaid'' sang ``Paid'' hoặc ''Overdue'' nếu quá hạn.
            \item Trong U3.3: Được gọi khi Nhân viên bấm "Xác nhận đã nhận tiền" để đổi trạng thái giao dịch sang "Đã thanh toán" (Success) rồi lưu xuống cơ sở dữ liệu.
        \end{itemize}
    \end{itemize}
    
    \item \textbf{calculateFee(): double (Class ParkingSession)}
    \begin{itemize}
        \item Mục đích: Tính toán chính xác số tiền phí mà khách hàng phải trả cho một phiên gửi xe.
        \item Mô tả logic: Tính toán số tiền (dưới dạng số thực double) dựa trên sự chênh lệch giữa \texttt{checkInTime} và \texttt{checkOutTime}, đối chiếu với \texttt{PricingPolicy} (bảng giá).
    \end{itemize}

    \item \textbf{isLocked(): boolean (Class Account)}
    \begin{itemize}
        \item Mục đích: Kiểm tra trạng thái hoạt động của tài khoản khách hàng (xem có đang bị khóa hay không).
        \item Mô tả logic: Phương thức này truy xuất giá trị của thuộc tính isLocked trong cơ sở dữ liệu của người dùng. Hệ thống sẽ gọi hàm này ở các bước kiểm tra điều kiện (Ví dụ: trước khi cho phép xe vào bãi, hoặc ngay khi khách hàng đăng nhập) để xem tài khoản có đang bị khóa do nợ quá hạn hay không.
    \end{itemize}
\end{itemize}

